% from main.tex
%%%%%%%%%%  scan idx 639  |  folio 626  %%%%%%%%%%
nombre d'intersection, qui peut être calculé comme coefficient du monôme \(n_1\cdots n_s\) dans le polynôme de Snapper \(\chi\bigl(L_1^{\otimes n_1}\otimes\cdots\otimes L_s^{\otimes n_s}\otimes\mathcal O_Y\bigr)\) (X 4.3.1).

\textbf{2.2.} Désormais, on suppose que \(k\) est algébriquement clos, que \(X\) est projectif, muni d'un faisceau inversible très ample \(H=\mathcal O_X(1)\) et que \(X\) est intègre, de dimension \(r\).

Soit \(L\) un faisceau inversible sur \(X\). L'entier \(\langle c_1(L)\cdot c_1(H)^{r-1}\rangle\) est appelé le \emph{degré} de \(L\) (relativement à \(H\)) et est noté \(d(L)\).

Dans la suite on considère un faisceau inversible \(L\) et les polynômes de Hilbert : \(\chi(L(n))=\sum_{i=0}^{r}a_i\binom{n+i}{i}\) et \(\chi(\mathcal O_X(n))=\sum_{i=0}^{r}e_i\binom{n+i}{i}\).

\textbf{Lemme 2.3.} \emph{On a} :
\begin{enumerate}[label=(\roman*)]
\item \emph{\(a_r=e_r=d(H)\).}
\item \emph{(Si \(r\ge1\)), \(a_{r-1}-e_{r-1}=d(L)\).}
\end{enumerate}

En effet, on obtient (i) et (ii) aussitôt en faisant \(m=0,1\) dans le polynôme suivant :
\[
\chi\bigl(L^{\otimes m}\otimes H^{\otimes n}\bigr)
=
\sum_{i+j\le r} f_{ij}\binom{m+i}{i}\binom{n+j}{j}.
\]

\textbf{Lemme 2.4.} \emph{Si \(d(L)\le0\) et \(h^0(L)\ne0\), alors \(L\simeq\mathcal O_X\); donc \(d(L)=0\) et \(h^0(L)=1\).}

En effet, comme \(X\) est intègre, il existe un diviseur \(Y\ge0\) tel que \(L\simeq\mathcal O_X(Y)\). Si \(Y>0\), on aurait \(0\ge\langle c(L)\cdot c_1(H)^{r-1}\rangle=\langle c_1(H)^{r-1}\cdot Y\rangle\ge1\).

\textbf{Lemme 2.5.} \emph{Si \(d(L)\le0\), alors \(L\) est un \((0,\ldots,0,a_r)\)-faisceau (avec \(r-1\) zéros).}

